\chapter*{Preface}
\addcontentsline{toc}{chapter}{Preface}

\section*{Why This Book?}

The landscape of biological sequencing technologies has expanded at a breathtaking pace. What began with Frederick Sanger's chain-termination method in 1977---capable of reading a few hundred bases at a time---has evolved into an ecosystem of technologies that can interrogate virtually every layer of biological information: the genome, the transcriptome, the epigenome, the three-dimensional organisation of chromatin, protein--nucleic acid interactions, and even the spatial arrangement of molecules within intact tissues.

For newcomers, this landscape can be overwhelming. Hundreds of protocol names, dozens of platforms, and a dizzying array of acronyms stand between a curious biologist and a productive sequencing experiment. This book aims to be the single reference that bridges that gap.

\section*{Who Is This Book For?}

This book assumes essentially no prior knowledge of molecular biology or bioinformatics. If you know that DNA is a molecule and that cells contain genes, you have enough background to start reading from \cref{ch:introduction}. We build every concept from the ground up:

\begin{itemize}[nosep]
  \item \textbf{Graduate students} entering a genomics or bioinformatics lab for the first time.
  \item \textbf{Wet-lab biologists} who want to understand what happens after they submit samples to a sequencing core.
  \item \textbf{Computational scientists} who want to understand the biology and chemistry behind the data they analyse.
  \item \textbf{Clinicians and translational researchers} exploring how sequencing technologies can be applied in clinical settings.
  \item \textbf{Industry professionals} evaluating sequencing platforms and applications for their organisations.
\end{itemize}

\section*{How to Read This Book}

The book is organised into seven parts:

\begin{description}[style=nextline]
  \item[Part~I: Foundations] covers the essential molecular biology, the history of sequencing, and the core principles shared by all modern sequencing technologies.
  \item[Part~II: Platforms and Technologies] provides a detailed survey of current sequencing instruments---short-read and long-read---and the universal concepts of library preparation.
  \item[Part~III: Genomics] covers DNA-level applications including whole-genome sequencing, exome sequencing, variant discovery, and metagenomics.
  \item[Part~IV: Transcriptomics] explores RNA sequencing in all its forms: bulk, single-cell, long-read, and spatial.
  \item[Part~V: Epigenomics] dives into DNA methylation, chromatin accessibility, histone modifications, and three-dimensional genome architecture.
  \item[Part~VI: Multiomics and Emerging Frontiers] covers the latest technologies that combine multiple measurement modalities and the cutting-edge methods pushing the field forward.
  \item[Part~VII: Bioinformatics and Data Analysis] walks through computational workflows, experimental design, data standards, and reproducibility---with concrete examples using both Nextflow and Snakemake workflow managers.
\end{description}

Each chapter can largely be read independently, though readers entirely new to the field will benefit from reading Parts~I and~II first. Throughout the text, coloured boxes highlight key concepts, practical tips, protocol summaries, and important caveats.

\section*{A Note on Pace of Change}

Sequencing technology evolves rapidly. We have endeavoured to cover the state of the art as of early 2026, but specific instrument models, software versions, and even company names may change. The fundamental principles---the chemistry, the biology, and the analytical strategies---evolve much more slowly and will remain relevant for years to come.

\vspace{1cm}
\noindent\textit{Happy sequencing.}
